% Part: proof-theory
% Chapter: cut-elimination
% Section: interpolation

\documentclass[../../../include/open-logic-section]{subfiles}

\begin{document}

\olfileid{pt}{cut}{itp}

\olsection{Maehara引理与 Craig插值定理}

插值定理由 William Craig（克雷格）提出，它表明：如果 $!A \Entails !B$，那么存在一个语句~$!C$（称为\emph{插值式}），使得 $!A \Entails !C$ 且 $!C \Entails !B$。重要的是，可以选取插值式~$!C$，使得出现在~$!C$ 中的所有常项符号和谓词符号均同时出现在~$!A$ 和~$!B$ 中。（我们假设这里不涉及函数符号。）若无此限制，$!A$ 和~$!B$ 本身就能平凡地作为插值式。

插值定理对经典一阶逻辑成立，并被称为一阶逻辑“最后被发现的重要性质”。
它在模型论和自动推理中具有重要应用。
其最初的动机与应用在于科学哲学，因为它可用于研究一个理论能否或如何借由某些初始概念来公理化。
它还蕴涵 Beth 可定义性定理，该定理指出，若一个理论能隐定义一个概念，则它也能显定义该概念。

与数理逻辑中的许多结果一样，插值定理既可以用模型论的方法证明，也可以用证明论的方法证明。
证明论方法的优势在于，不仅能证明插值式的存在，还能提供一个算法，从 $!A \Entails !B$ 的推导（例如，~\Log{G3c} 中 $!A \Sequent !B$ 的推导）计算出该插值式。

设 $\Lang L(!A)$ 为 $!A$ 中的常项符号与谓词符号的集合，并设 $\Lang L(\Gamma) = \bigcup \Setabs{\Lang L(!A)}{!A \in
\Gamma}$。
本节始终假设所处理的公式不含函数符号。

\begin{thm}[Craig 插值定理]\ollabel{thm:interpolation}
设 $!A$ 的语言为 $\Lang L(!A)$，$!B$ 的语言为 $\Lang L(!B)$。
如果 $\Log{G3c} \Proves !A \Sequent !B$，则存在公式~$!C$，其语言为 $\Lang L(!C) \subseteq \Lang
L(!A) \cap \Lang L(!B)$，使得 $\Log{G3c} \Proves !A \Sequent !C$ 且 $\Log{G3c} \Proves !C \Sequent !B$。
\end{thm}

我们证明插值定理的一个推广，以便利用~\Log{G3c} 中推导的结构。
为此，设想将相继式的前件和后件各分为两部分。
我们记作 $\maeh{\Gamma_1}{\Gamma_2} \Sequent
\maeh{\Delta_1}{\Delta_2}$，表示 $\Gamma_1$ 和 $\Delta_1$ 属于第一部分，而 $\Gamma_2$ 和 $\Delta_2$ 属于第二部分。
于是，插值定理便是下述引理的直接推论。

\begin{lem}[Maehara 引理]\ollabel{lem:maehara}
如果 $\Log{G3c}
\Proves \maeh{\Gamma_1}{\Gamma_2} \Sequent \maeh{\Delta_1}{\Delta_2}$，则存在公式~$!C$，其语言满足 $\Lang L(!C)
\subseteq \Lang L(\Gamma_1,
\Delta_1) \cap \Lang L(\Gamma_2, \Delta_2)$，使得 $\Log{G3c}
\Proves \Gamma_1 \Sequent \Delta_1, !C$ 且 $\Log{G3c} \Proves !C,
\Gamma_2 \Sequent \Delta_2$。
\end{lem}

\begin{proof}
  假设 $\pi \Proves \maeh{\Gamma_1}{\Gamma_2} \Sequent
  \maeh{\Delta_1}{\Delta_2}$。我们对 $n=\pheight{\pi}$ 进行归纳来证明该断言。

  如果 $n=0$，那么 $\maeh{\Gamma_1}{\Gamma_2} \Sequent
  \maeh{\Delta_1}{\Delta_2}$ 是一个公理，并且某个原子公式~$!A$ 同时出现在 $\Gamma_1, \Gamma_2$ 和 $\Delta_1,
  \Delta_2$ 中。根据 $!A$~属于左侧的第一部分还是第二部分，以及属于右侧的第一部分还是第二部分，分为四种情况。
  \begin{enumerate}
    \item $!A \in \Gamma_1$ 且 $!A \in \Delta_1$。我们需要找到一个~$!C$，使得 $\Gamma_1 \Sequent \Delta_1, !C$ 和 $!C, \Gamma_2
    \Sequent \Delta_2$ 都是可推导的。因为 $!A$ 同时出现在左侧和右侧，无论我们如何选取~$!C$，第一个相继式都已经是公理。为了保证 $!C, \Gamma_2
    \Sequent \Delta_2$ 是可推导的，我们唯一的选择是令 $!C \ident \lfalse$。
    \item $!A \in \Gamma_1$ 且 $!A \in \Delta_2$。那么 $\Gamma_1
    \Sequent \Delta_1, !A$ 和 $!A, \Gamma_2 \Sequent \Delta_2$ 都是公理，因此我们可以选取 $!C \ident !A$。
    \item $!A \in \Gamma_2$ 且 $!A \in \Delta_1$。那么 $!A, \Gamma_1
    \Sequent \Delta_1$ 和 $\Gamma_2 \Sequent \Delta_2, !A$ 将是公理，但 $!A$ 出现在了错误的一侧。因此 $!A$ 不是插值。然而，使用 $\RightR{\lnot}$ 和 $\LeftR{\lnot}$ 规则，我们可以推导出 $\Gamma_1
    \Sequent \Delta_1, \lnot!A$ 和 $\lnot !A, \Gamma_2 \Sequent
    \Delta_2$。
    \item $!A \in \Gamma_2$ 且 $!A \in \Delta_2$。练习。
  \end{enumerate}
  因 $\lfalse \in \Gamma_1$ 或 $\lfalse \in \Gamma_2$ 而使得 $\maeh{\Gamma_1}{\Gamma_2} \Sequent
  \maeh{\Delta_1}{\Delta_2}$ 是公理的情况，也留作练习。

  如果 $n>0$，则 $\pi$ 以一条逻辑推理规则结束。我们必须根据使用了哪条规则以及主公式属于哪个部分来区分情况。在每种情况下，我们都可以使用归纳假设：该引理对所有满足 $\pheight{\pi_1} < n$ 的推导~$\pi_1$ 成立，且无论其末尾相继式如何被分割为两部分。让我们考虑一些具体的情况。
  \begin{enumerate}
    \item $\pi$ 以 $\RightR{\lnot}$ 结束，主公式为~$\lnot !A$。我们有两种情况：根据 $\lnot !A$ 属于第一部分还是第二部分，$\pi$ 要么结束于
    \[
    \AxiomC{}
    \RightLabel{$\pi_1$}
    \Deduce$\maeh{!A, \Gamma_1}{\Gamma_2} \fCenter
      \maeh{\Delta_1}{\Delta_2}$
    \RightLabel{\RightR{\lnot}}
    \UnaryInf$\maeh{\Gamma_1}{\Gamma_2} \fCenter
      \maeh{\Delta_1, \lnot !A}{\Delta_2}$
    \DisplayProof
    \]
    要么
    \[
    \AxiomC{}
    \RightLabel{$\pi_1$}
    \Deduce$\maeh{\Gamma_1}{!A, \Gamma_2} \fCenter
      \maeh{\Delta_1}{\Delta_2}$
    \RightLabel{\RightR{\lnot}}
    \UnaryInf$\maeh{\Gamma_1}{\Gamma_2} \fCenter
      \maeh{\Delta_1}{\Delta_2, \lnot !A}$
    \DisplayProof
    \]
    注意，如果主公式~$\lnot !A$ 属于第一部分，我们也将前提中的活动公式~$!A$ 分配给了第一部分。这是我们可以做出的选择。如果我们把活动公式分配给另一部分，归纳假设同样适用，但那样我们将无法找到一个插值（练习：尝试一下）。

    首先考虑第一种情形，即 $\lnot !A$~属于第一部分。根据归纳假设，存在~$\pi_1$ 的末尾相继式的一个插值，即一个公式~$!C_1$，使得
    \begin{align*}
    !A, \Gamma_1 & \Sequent \Delta_1, !C_1 &&与 &
    !C_1, \Gamma_2 & \Sequent \Delta_2
    \intertext{都是可证的。我们需要的是一个公式~$!C$，使得以下两个相继式都是可推导的：}
    \Gamma_1 & \Sequent \Delta_1, \lnot !A, !C &&与 &
    !C, \Gamma_2 & \Sequent \Delta_2
    \end{align*}
    显然，我们可以直接取 $!C \ident !C_1$：归纳假设已经告诉我们右侧的相继式是可推导的，而左侧的相继式则可以通过对归纳假设得到的相继式应用~$\RightR{\lnot}$ 规则得到。归纳假设还告诉我们 $\Lang L(!C_1) \subseteq \Lang L(!A, \Gamma_1,
    \Delta_1) \cap \Lang L(\Gamma_2, \Delta_2)$。由于 $\Lang L(\lnot
    !A) = \Lang L(!A)$ 且 $\Lang L(!C) = \Lang L(!C_1)$，我们有 $\Lang L(!C_1) \subseteq \Lang L(\Gamma_1, \Delta_1, \lnot !A) \cap \Lang
    L(\Gamma_2, \Delta_2)$。

    在第二种情形中，情况略有不同，因为 $\lnot !A$~属于第二部分。我们将前提中的活动公式也分配给第二部分，并应用归纳假设可知
    \begin{align*}
    \Gamma_1 & \Sequent \Delta_1, !C_1 &&与 &
    !C_1, !A, \Gamma_2 & \Sequent \Delta_2
    \intertext{是可证的。如果我们再次对第二个相继式应用 \RightR{\lnot} 规则，就得到}
    \Gamma_1 & \Sequent \Delta_1, !C_1 &&与 &
    !C_1, \Gamma_2 & \Sequent \Delta_2, \lnot !A
    \end{align*}的推导。
    这些正是在 $!C_1$ 是插值时我们希望可推导的相继式；我们再次令 $!C \ident !C_1$。和前面一样，有 $\Lang
    L(!C_1) \subseteq \Lang L(\Gamma_1, \Delta_1) \cap \Lang L(\Gamma_2,
    \Delta_2, \lnot !A)$。
    \item $\pi$ 以 $\RightR{\lif}$ 结束，主公式为~$!A \lif !B$。我们同样有两种情况，取决于 $!A \lif !B$ belongs to the first or second part. The !!{proof}~$$，$\pi$ 要么结束于
    \[
    \AxiomC{}
    \RightLabel{$\pi_1$}
    \Deduce$\maeh{!A, \Gamma_1}{\Gamma_2} \fCenter
      \maeh{\Delta_1, !B}{\Delta_2}$
    \RightLabel{\RightR{\lif}}
    \UnaryInf$\maeh{\Gamma_1}{\Gamma_2} \fCenter
      \maeh{\Delta_1, !A \lif !B}{\Delta_2}$
    \DisplayProof
    \]
    要么
    \[
    \AxiomC{}
    \RightLabel{$\pi_1$}
    \Deduce$\maeh{\Gamma_1}{!A, \Gamma_2} \fCenter
      \maeh{\Delta_1}{\Delta_2, !B}$
    \RightLabel{\RightR{\lif}}
    \UnaryInf$\maeh{\Gamma_1}{\Gamma_2} \fCenter
      \maeh{\Delta_1}{\Delta_2, !A \lif !B}$
    \DisplayProof
    \]
    同样地，我们在前提中将活动公式分配给了与结论中主公式相同的部分。在第一种情形中，我们应用归纳假设，得到
    \begin{align*}
    !A, \Gamma_1 & \Sequent \Delta_1, !B, !C_1 &&与 & !C_1, \Gamma_2 & \Sequent \Delta_2
    \intertext{的推导。左侧的相继式是应用 \RightR{\lif} 规则的合适前提。因此我们有以下可推导的相继式：}
    \Gamma_1 & \Sequent \Delta_1, !A \lif !B, !C_1 &&与 & !C_1, \Gamma_2 & \Sequent \Delta_2
    \end{align*}
    这意味着 $!C_1$ 也是~$\pi$ 的末尾相继式的一个插值，我们可以再次取 $!C \ident !C_1$。

    另一种情形的处理方式相同。
    \item $\pi$ 以 $\LeftR{\lif}$ 结束，主公式为~$!A \lif !B$。如果 $!A \lif !B$ 属于第一部分，推导~$\pi$ 结束于
    \[
    \AxiomC{}
    \RightLabel{$\pi_1$}
    \Deduce$\maeh{\Gamma_1}{\Gamma_2} \fCenter
      \maeh{\Delta_1, !A}{\Delta_2}$
    \AxiomC{}
    \RightLabel{$\pi_2$}
    \Deduce$\maeh{!B, \Gamma_1}{\Gamma_2} \fCenter
      \maeh{\Delta_1}{\Delta_2}$
    \BinaryInf$\maeh{!A \lif !B, \Gamma_1}{\Gamma_2} \fCenter
      \maeh{\Delta_1}{\Delta_2}$
    \DisplayProof
    \]
    现在归纳假设给我们四个可推导的相继式，其中两个对应于左前提的插值 $!C_1$，另两个对应于右前提的插值 $!C_2$：
    \begin{align*}
    \Gamma_1 & \Sequent \Delta_1, !A, !C_1 &&\text{(1)} &
    !C_1, \Gamma_2 & \Sequent \Delta_2 && \text{(2)}\\
    !B, \Gamma_1 & \Sequent \Delta_1, !C_2 &&\text{(3)} &
    !C_2, \Gamma_2 & \Sequent \Delta_2 && \text{(4)}
    \end{align*}
    如果我们应用弱化规则，在 (1) 的右侧添加 $!C_2$，在 (3) 的右侧添加 $!C_1$，相继式 (1) 和 (3) 就可以作为~\LeftR{\lif} 规则的前提：
    \[
    \AxiomC{}
    \Deduce$\Gamma_1 \fCenter \Delta_1, !A, !C_1, !C_2$
    \AxiomC{}
    \Deduce$!B, \Gamma_1 \fCenter \Delta_1, !C_1, !C_2$
    \RightLabel{\LeftR{\lif}}
    \BinaryInf$!A \lif !B, \Gamma_1 \fCenter \Delta_1, !C_1, !C_2$
    \DisplayProof
    \]
    应用 $\RightR{\lor}$ 规则，我们得到相继式
    \[
    !A \lif !B, \Gamma_1 \fCenter \Delta_1, !C_1 \lor !C_2
    \]
    这表明在这种情况下，$!C_1 \lor !C_2$ 可能是此时的插值。确实，我们可以从剩下的相继式 (2) 和 (4) 推导出另一个所需的相继式：
    \[
    \AxiomC{}
    \Deduce$!C_1, \Gamma_2 \fCenter \Delta_2$
    \AxiomC{}
    \Deduce$!C_2, \Gamma_2 \fCenter \Delta_2$
    \RightLabel{\LeftR{\lor}}
    \BinaryInf$!C_1 \lor !C_2, \Gamma_2 \fCenter \Delta_2$
    \DisplayProof
    \]
    接下来需要验证 $!C_1 \lor !C_2$ 属于正确的语言。归纳假设告诉我们
    \begin{align*}
    \Lang L(!C_1) & \subseteq
    \Lang L(\Gamma_1, \Delta_1, !A) \cap \Lang L(\Gamma_2, \Delta_2) &与\\
    \Lang L(!C_2) & \subseteq
    \Lang L(!B, \Gamma_1, \Delta_1) \cap \Lang L(\Gamma_2, \Delta_2)
    \intertext{由于}
    \Lang L(!A) & \subseteq \Lang L(!A \lif !B) &与\\
    \Lang L(!B) & \subseteq \Lang L(!A \lif !B)
    \intertext{我们既有}
    \Lang L(!C_1) & \subseteq
    \Lang L(\Gamma_1, \Delta_1, !A \lif !B) \cap \Lang L(\Gamma_2, \Delta_2) &与\\
    \Lang L(!C_2) & \subseteq
    \Lang L(\Gamma_1, \Delta_1, !A \lif !B) \cap \Lang L(\Gamma_2, \Delta_2)
    \intertext{因此，$\Lang L (!C_1 \lor !C_2) = {}$}
    \Lang L(!C_1) \cup \Lang L(!C_2) & \subseteq
    \Lang L(\Gamma_1, \Delta_1, !A \lif !B) \cap \Lang L(\Gamma_2, \Delta_2),
    \end{align*}
    这符合要求。

    在 $!A \lif !B$ 属于第二部分的情况下，情况有所不同，但处理方式类似。归纳假设再次给我们四个可推导的相继式：
    \begin{align*}
    \Gamma_1 & \Sequent \Delta_1, !C_1 &&\text{(1)} &
    !C_1, \Gamma_2 & \Sequent \Delta_2, !A && \text{(2)}\\
    \Gamma_1 & \Sequent \Delta_1, !C_2 &&\text{(3)} &
    !C_2, !B, \Gamma_2 & \Sequent \Delta_2 && \text{(4)}
    \end{align*}
    现在，相继式 (2) 和 (4)——加上一些通过弱化添加的公式——可以作为~\LeftR{\lif} 规则的前提：
    \[
    \AxiomC{}
    \Deduce$!C_1, !C_2, \Gamma_2 \fCenter \Delta_2, !A$
    \AxiomC{}
    \Deduce$!B, !C_1, !C_2, \Gamma_2 \fCenter \Delta_2$
    \RightLabel{\LeftR{\lif}}
    \BinaryInf$!A \lif !B, !C_1, !C_2, \Gamma_2 \fCenter \Delta_2$
    \DisplayProof
    \]
    我们再次希望将其转化为前件包含一个由 $!C_1$ 和 $!C_2$ 构成的公式的相继式（因为现在处理的是第二部分）。应用 $\LeftR{\land}$ 规则即可；我们应将 $!C \ident
    (!C_1 \land !C_2)$ 取作插值。碰巧，我们可以用剩下的相继式 (1) 和 (3) 推导出对应另一部分的相继式：
    \[
    \AxiomC{}
    \Deduce$\Gamma_1 \fCenter \Delta_1, !C_1$
    \AxiomC{}
    \Deduce$\Gamma_1 \fCenter \Delta_1, !C_2$
    \RightLabel{\RightR{\land}}
    \BinaryInf$\Gamma_1 \fCenter \Delta_1, !C_1 \land !C_2$
    \DisplayProof
    \]
    和前面一样，我们有 $\Lang L(!C_1 \land !C_2) \subseteq \Lang
    L(\Gamma_1, \Delta_1) \cap \Lang L(\Gamma_2, \Delta_2, !A \lif
    !B)$。

    \item $\pi$ 以 $\RightR{\lforall}$ 结束，主公式为~$\lforall[x][!A(x)]$：
    \[
    \AxiomC{}
    \RightLabel{$\pi_1$}
    \Deduce$\maeh{\Gamma_1}{\Gamma_2} \fCenter
    \maeh{\Delta_1, !A(c)}{\Delta_2}$
    \RightLabel{\RightR{\lforall}}
    \UnaryInf$\maeh{\Gamma_1}{\Gamma_2} \fCenter
      \maeh{\Delta_1, \lforall[x][!A(x)]}{\Delta_2}$
    \DisplayProof
    \]
    或
    \[
    \AxiomC{}
    \RightLabel{$\pi_1$}
    \Deduce$\maeh{\Gamma_1}{\Gamma_2} \fCenter
      \maeh{\Delta_1}{\Delta_2, !A(c)}$
    \RightLabel{\RightR{\lforall}}
    \UnaryInf$\maeh{\Gamma_1}{\Gamma_2} \fCenter
      \maeh{\Delta_1}{\Delta_2, \lforall[x][!A(x)]}$
    \DisplayProof
    \]
    在第一种情形中，归纳假设给出了 $\Gamma_1 \Sequent \Delta_1, !A(c), !C_1$ 和 $!C_1, \Gamma_2
    \Sequent \Delta_2$ 的推导。我们可以对第一个相继式应用 $\RightR{\lforall}$ 规则得到 $\Gamma_1 \Sequent \Delta_1, \lforall[x][!A(x)], !C_1$。
    （特征变元条件满足，因为它在~$\pi$ 末尾的 \RightR{\lforall} 规则中已经满足。）因此，如果 $\Lang
    L(!C_1) \subseteq \Lang L(\Gamma_1, \Delta_1, \lforall[x][!A(x)])
    \cap \Lang L(\Gamma_2, \Delta_2)$ 这一条件满足，那么 $!C_1$ 就是一个插值。根据归纳假设，我们有 $\Lang L(!C_1) \subseteq \Lang L(\Gamma_1,
    \Delta_1, !A(c)) \cap \Lang L(\Gamma_2, \Delta_2)$。所以我们需要担心~$c$——$c$ 可能出现在~$!C_1$ 中吗？不可能。如果 $c$ 出现在 $!C_1$ 中，我们将同时有 $c \in \Lang L(\Gamma_1, \Delta_1,
    !A(c))$（这是真的）以及 $c \in \Lang L(\Gamma_2,
    \Delta_2)$。但这样一来，$c$ 就会出现在~$\pi$ 中 $\RightR{\lforall}$ 推理的结论里，违反了特征变元条件。

    另一种情形类似。

    \item $\pi$ 以 $\RightR{\lexists}$ 结束，主公式为~$\lexists[x][!A(x)]$。我们同样有两种情形。我们考虑 $\lexists[x][!A(x)]$ 属于第一部分的情形，另一种情形留作练习。

    在第一种情形中，$\pi$ 结束于
    \[
    \AxiomC{}
    \RightLabel{$\pi_1$}
    \Deduce$\maeh{\Gamma_1}{\Gamma_2} \fCenter
      \maeh{\Delta_1, \lexists[x][!A(x)], !A(t)}{\Delta_2}$
    \RightLabel{\RightR{\lexists}}
    \UnaryInf$\maeh{\Gamma_1}{\Gamma_2} \fCenter
      \maeh{\Delta_1, \lexists[x][!A(x)]}{\Delta_2}$
    \DisplayProof
    \]
    归纳假设给出了 $\Gamma_1 \Sequent
    \Delta_1, \lexists[x][!A(x)], !A(t), !C_1$ 和 $!C_1, \Gamma_2
    \Sequent \Delta_2$ 的推导。我们可以对第一个相继式应用 $\RightR{\lexists}$ 规则得到 $\Gamma_1 \Sequent \Delta_1, \lexists[x][!A(x)], !C_1$。
    因此，如果条件 $\Lang L(!C_1) \subseteq \Lang L(\Gamma_1, \Delta_1,
    \lexists[x][!A(x)]) \cap \Lang L(\Gamma_2, \Delta_2)$ 满足，$!C_1$ 将再次是一个插值。这里归纳假设告诉我们 $\Lang L(!C_1)
    \subseteq \Lang L(\Gamma_1, \Delta_1, \lexists[x][!A(x)], !A(t))
    \cap \Lang L(\Gamma_2, \Delta_2)$. Now remember that we assumed
    that there are no !!{function}s. That means the term~$t$ can only
    be !!a{constant}, i.e., $t \ident b$。如果 $b \notin \Lang L(!C_1)$
    或者 $b \in \Lang L(\Gamma_1, \Delta_1, \lexists[x][!A(x)]) \cap
    \Lang L(\Gamma_2, \Delta_2)$，那么没什么可担心的：我们可以取 $!C_1$ 作为插值。但如果 $b \in \Lang
    L(!C_1)$ 且 $b \notin \Lang L(\Gamma_1, \Delta_1,
    \lexists[x][!A(x)]) \cap \Lang L(\Gamma_2, \Delta_2)$? First of
    all, by replacing all occurrences of~$$，将~$!C_1$ 中的 $b$ 替换为变元~$y$，我们可以写成 $!C_1 \ident \Subst{!C_1'}{b}{y}$
    其中 $!C_1'$~不含~$b$。此外，要么 $b \notin
    \Lang L(\Gamma_1, \Delta_1, \lexists[x][!A(x)])$，要么 $b \notin
    \Lang L(\Gamma_2, \Delta_2)$。我们可以分别取 $!C \ident
    \lforall[y][!C_1'(y)]$ 或 $!C \ident \lexists[y][!C_1'(y)]$。
    在第一种情况下，我们有：
    \[
    \AxiomC{}
    \Deduce$\Gamma_1 \fCenter
      \Delta_1, \lexists[x][!A(x)], !A(b), !C_1'(b)$
    \RightLabel{\RightR{\lexists}}
    \UnaryInf$\Gamma_1 \fCenter \Delta_1, \lexists[x][!A(x)], !C_1'(b)$
    \RightLabel{\RightR{\lforall}}
    \UnaryInf$\Gamma_1 \fCenter
      \Delta_1, \lexists[x][!A(x)], \lforall[y][!C_1'(y)]$
    \DisplayProof
    \]
    和
    \[
    \AxiomC{}
    \Deduce$!C_1'(b), \lforall[y][!C_1'(y)], \Gamma_2 \fCenter \Delta_2$
    \RightLabel{\LeftR{\lforall}}
    \UnaryInf$\lforall[y][!C_1'(y)], \Gamma_2 \fCenter \Delta_2$
    \DisplayProof
    \]
    顶端相继式的推导来自归纳假设（在第二个推导中，我们在前件应用了对 $\lforall[y][!C_1'(y)]$ 的弱化可容许性）。第一个推导中 \RightR{\lforall} 规则的特征变元条件得到满足，因为 $b \notin \Lang L(\Gamma_1, \Delta_1,
    \lexists[x][!A(x)])$，且 $!C_1'$~的定义保证了它不包含~$b$。

    在第二种情况下，我们有（基于归纳假设及对 $\lexists[y][!C_1'(y)]$ 的弱化可容许性）：
    \[
    \AxiomC{}
    \Deduce$\Gamma_1 \fCenter
      \Delta_1, \lexists[x][!A(x)], !A(b), \lexists[y][!C_1'(y)], !C_1'(b)$
    \RightLabel{\RightR{\lexists}}
    \UnaryInf$\Gamma_1 \fCenter
      \Delta_1, \lexists[x][!A(x)], \lexists[y][!C_1'(y)], !C_1'(b)$
    \RightLabel{\RightR{\lexists}}
    \UnaryInf$\Gamma_1 \fCenter
      \Delta_1, \lexists[x][!A(x)], \lexists[y][!C_1'(y)]$
    \DisplayProof
    \]
    和
    \[
    \AxiomC{}
    \Deduce$!C_1'(b), \Gamma_2 \fCenter \Delta_2$
    \RightLabel{\LeftR{\lexists}}
    \UnaryInf$\lexists[y][!C_1'(y)], \Gamma_2 \fCenter \Delta_2$
    \DisplayProof
    \]
    第二个推导中 \LeftR{\lexists} 规则的特征变元条件得到满足，因为 $b \notin \Lang L(\Gamma_2,
    \Delta_2)$，且 $!C_1'$ 的定义保证了它不包含~$b$。
  \end{enumerate}
  其他情况类似；我们将其留作练习。
\end{proof}

\begin{prob}
假设我们有一个联结词 $\lnand$ 及其规则
\[
\Axiom$\Gamma \fCenter \Delta, !A$
\Axiom$\Gamma \fCenter \Delta, !B$
\RightLabel{\LeftR{\lnand}}
\BinaryInf$!A \lnand !B, \Gamma \fCenter \Delta$
\DisplayProof
\quad
\Axiom$!A, !B, \Gamma \fCenter \Delta$
\RightLabel{\RightR{\lnand}}
\UnaryInf$\Gamma \fCenter \Delta, !A \lnand !B$
\DisplayProof
\]
通过处理 \olref{lem:maehara} 证明中 $\pi$ 以这些规则之一结尾的情形，证明 Maehara 引理仍然成立。
\end{prob}

确立了 Maehara 引理之后，我们可以用它来证明 Craig 插值定理。

\begin{proof}[Proof of \olref{thm:interpolation}]
  假设 $!A \Sequent !B$ 是可推导的。将 \olref{lem:maehara} 应用于 $\maeh{!A}{\ } \Sequent \maeh{\ }{!B}$ 以得到
  插值式~$!C$。则有 $\Proves !A \Sequent !C$ 和 $\Proves !C
  \Sequent !B$。
\end{proof}

假设 $\Gamma$ 是一个语言~$\Lang L$ 中的理论（语句集合），该语言包含一个 $n$ 元谓词~$R$。我们说 $\Gamma$ \emph{隐式定义}了~$R$，当且仅当
\begin{align*}
\Gamma, \Gamma'
& \Entails \lforall[x_1][\dots\lforall[x_n][(R(x_1,\dots,x_n) \liff
R'(x_1, \dots, x_n))]],
\intertext{其中 $\Gamma'$ 是将 $\Gamma$ 中所有 $R$ 的出现替换为不在 $\Lang L$ 中的 $R'$ 后得到的。$\Gamma$ \emph{显式定义}了~$R$，当且仅当存在某个不包含~$R$ 的公式~$!A(x_1, \dots, x_n)$ 使得}
\Gamma &\Entails
\lforall[x_1][\dots\lforall[x_n][(R(x_1,\dots,x_n) \liff !A(x_1,
\dots, x_n))]].
\end{align*}
显然，如果 $\Gamma$ 显式定义了 $R$，那么它也就隐式定义了~$R$。反之则并不显然，但依然成立：

\begin{lem}[Beth 可定义性定理]
  如果 $\Gamma$ 隐式定义了~$R$，那么它也显式定义了~$R$。
\end{lem}

\begin{proof}
  假设 $\Gamma$ 隐式定义了~$R$。令 $R'$、$\Gamma'$ 如上所述，并令 $\Lang L' = \Lang L(\Gamma') = \Lang L \setminus \{R\}
  \cup \{R'\}$。由假设，我们有
  \begin{align*}
    \Gamma, \Gamma' &
    \Sequent \lforall[x_1][\dots\lforall[x_n][(R(x_1,\dots,x_n)
    \liff R'(x_1, \dots, x_n))]]
  \intertext{令 $c_1$, \dots,~$c_n$ 为不在 $\Gamma$ 中出现的常项符号。由可逆性引理，我们得到}
    \Gamma, \Gamma', R(c_1,\dots,c_n) & \Sequent  R'(c_1, \dots, c_n)
  \intertext{将 Maehara 引理应用于}
    \maeh{\Gamma, R(c_1,\dots,c_n)}{\Gamma'} & \Sequent
      \maeh{\ }{R'(c_1, \dots, c_n)}
  \intertext{以找到一个公式~$!A \ident !A(x_1, \dots, x_n)$，使得}
    \Gamma, R(c_1,\dots,c_n) & \Sequent !A(c_1, \dots, c_n) \\
    !A(c_1, \dots, c_n), \Gamma' & \Sequent R'(c_1, \dots, c_n)
  \intertext{有推导。由于 $\Lang L(!A) = \Lang L_1 \cap
  \Lang L_2 = \Lang L(\Gamma) \setminus \{R\}$，$!A$ 不包含~$R$。
  将第二个相继式的推导中所有的 $R'$ 替换为~$R$，得到}
    !A(c_1, \dots, c_n), \Gamma & \Sequent R(c_1, \dots, c_n)
  \end{align*}
  应用 \RightR{\lif} 和 \RightR{\lforall} 规则，得到一个推导：
  \[
    \Gamma \Sequent \lforall[x_1][\dots\lforall[x_n][(R(x_1,\dots,x_n) \liff !A(x_1, \dots, x_n))]].
  \]
  这表明 $!A(x_1, \dots, x_n)$ 在~$\Gamma$ 中显式定义了 $!R$。
\end{proof}

第三个与插值性等价的结果（并且同样可以用 Maehara 引理证明）是下面的定理：如果两个一致理论~$\Gamma_1$、$\Gamma_2$ 包含一个语言~$\Lang L(\Gamma) \subseteq \Lang L(\Gamma_1)
\cap \Lang L(\Gamma_2)$ 上的完备理论 $\Gamma$，那么 $\Gamma_1 \cup \Gamma_2$ 是一致的。

回忆一下，完备理论指的是对其语言中的每个语句~$!A$，要么 $\Gamma \Proves !A$，要么 $\Gamma \Proves \lnot !A$。由于 $\Gamma_1$ 和~$\Gamma_2$ 是 $\Gamma$ 的一致扩张，$\Gamma$ 必须是一致的。由此可知，如果 $!A$ 是语言 $\Lang L(\Gamma_1) \cap \Lang
L(\Gamma_2)$ 中的一个语句，我们不可能有 $\Gamma_1 \Entails !A$ 且 $\Gamma_2
\Entails \lnot !A$。为了看出这一点，假设存在这样的~$!A$。如果 $\Gamma_1 \Entails !A$，那么 $\Gamma \Entails !A$。否则，由 $\Gamma$ 的完备性，$\Gamma \Entails \lnot !A$，并且由于 $\Gamma \subseteq \Gamma_1$，有 $\Gamma_1 \Entails \lnot !A$，这将导致 $\Gamma_1$ 不一致。因此 $\Gamma \Entails !A$，并且由于 $\Gamma
\subseteq \Gamma_2$，有 $\Gamma_2 \Entails !A$。所以 $\Gamma \Entails/
\lnot !A$，否则 $\Gamma_2$ 就会不一致。

在共同语言~$\Lang L(\Gamma_1) \cap \Lang L(\Gamma_2)$ 中不存在满足 $\Gamma_1
\Entails !A$ 和 $\Gamma_2 \Entails \lnot !A$ 的语句~$!A$，这正是证明所需的条件。下面我们用证明论方法，借助 Maehara 引理给出证明。

\begin{thm}[Robinson 联合一致性定理]
假设 $\Gamma_1$、$\Gamma_2$ 是一致的理论，并且不存在语言 $\Lang L(\Gamma_1) \cap \Lang L(\Gamma_2)$ 中的语句~$!A$ 使得 $\Gamma_1 \Entails !A$ 且 $\Gamma_2 \Entails \lnot !A$。那么 $\Gamma_1 \cup \Gamma_2$ 是一致的。
\end{thm}

\begin{proof}
  反设 $\Gamma_1 \cup \Gamma_2$ 不一致。那么相继式 $\Gamma_1, \Gamma_2 \Sequent \ $ 有一个推导。将 Maehara 引理应用于 $\maeh{\Gamma_1}{\Gamma_2}
  \Sequent \maeh{\ }{\ }$，得到一个语言~$\Lang L(\Gamma_1) \cap \Lang L(\Gamma_2)$ 中的语句~$!A$，使得
  $\Proves \Gamma_1 \Sequent !A$ 和 $\Proves !A, \Gamma_2 \Sequent
  \quad$。由后者可得 $\Proves \Gamma_2 \Sequent \lnot !A$。
  但这与我们假设不存在这样的语句相矛盾。
\end{proof}

以上我们默认了理论 $\Gamma$、$\Gamma_1$、$\Gamma_2$ 是有限的。根据紧致性定理，这些结果对无限理论也成立。

\end{document}